\documentclass[10.5pt,a4paper]{article}
\usepackage[utf8]{inputenc}
\usepackage[T1]{fontenc}
\IfFileExists{brazil.ldf}{\usepackage[brazil]{babel}}{}
\IfFileExists{charter.sty}{\usepackage{charter}}{\usepackage{lmodern}}
\IfFileExists{charter.sty}{\renewcommand{\ttdefault}{pcr}}{}
\usepackage[a4paper,top=2cm,bottom=2cm,left=2.2cm,right=2.2cm]{geometry}
\usepackage{microtype}
\usepackage{amsmath,amssymb}
\usepackage{booktabs}
\usepackage{longtable}
\usepackage{tabularx}
\usepackage{multirow}
\usepackage{enumitem}
\usepackage[dvipsnames,table]{xcolor}
\usepackage{titlesec}
\usepackage{fancyhdr}
\usepackage{listings}
\usepackage[most]{tcolorbox}
\usepackage{tikz}
\usetikzlibrary{positioning,arrows.meta,calc}
\usepackage{pgfplots}
\pgfplotsset{compat=1.18}
\usepackage{xstring}
\usepackage[font=small,labelfont=bf,labelsep=colon]{caption}
\captionsetup[figure]{name=Figura}
\captionsetup[table]{name=Tabela}
\usepackage{xurl}
\usepackage[colorlinks=true,linkcolor=azul,urlcolor=azul,citecolor=azul]{hyperref}

% Gerado por scripts/gerar_numeros_latex_2t.py. NAO editar a mao.
\newcommand{\horacorte}{sábado, 24/10/2026, às 20h00}
\newcommand{\statusprevisao}{preliminar}
\newcommand{\pFlavio}{51,3}
\newcommand{\pLula}{48,7}
\newcommand{\pAbst}{21,1}
\newcommand{\pBrancos}{1,8}
\newcommand{\pNulos}{2,9}
\newcommand{\nPesq}{26}
\newcommand{\nInst}{14}
\newcommand{\nInstUmT}{5}
\newcommand{\nHist}{111}
\newcommand{\dataUltimaPesq}{21/10}
\newcommand{\sigmaZero}{1,5}
\newcommand{\medErroUmT}{2,42}
\newcommand{\maeVenc}{0,65}
\newcommand{\maeMZero}{0,92}
\newcommand{\deltaVenc}{0,27}
\newcommand{\seVenc}{0,14}
\newcommand{\loeoVenc}{0,48}
\newcommand{\loeoMZero}{0,77}
\newcommand{\lulaUmT}{45,16}
\newcommand{\flavioUmT}{47,03}
\newcommand{\tercUmT}{7,81}
\newcommand{\flavioDoisCand}{51,0}
\newcommand{\lulaDoisCand}{49,0}
\newcommand{\tauMediana}{0,45}
\newcommand{\tauMedia}{0,61}
\newcommand{\tauImplFlavio}{55}
\newcommand{\tauImplLula}{45}
\newcommand{\priorMedianaFlavio}{51,4}
\newcommand{\priorMediaFlavio}{50,1}
\newcommand{\sensMZero}{51,3}
\newcommand{\sensMUm}{51,1}
\newcommand{\sensPersist}{51,7}
\newcommand{\sensMin}{50,1}
\newcommand{\sensMax}{51,8}
\newcommand{\nSens}{11}
\newcommand{\abstUmT}{21,08}
\newcommand{\brancosUmT}{1,84}
\newcommand{\nulosUmT}{2,93}
\newcommand{\govShareVinteSeis}{7,9}
\newcommand{\govNVinteSeis}{6}
\newcommand{\govShareHist}{41,6}
\newcommand{\deltaAbstCom}{1,72}
\newcommand{\deltaAbstSem}{1,60}
\newcommand{\pesoMax}{13,4}
\newcommand{\instPesoMax}{Quaest}
\newcommand{\pesoMin}{3,5}
\newcommand{\instPesoMin}{CNT/MDA}
\newcommand{\rotulosPesos}{{CNT/MDA},{Real Time Big Data},{Datafolha},{Palver},{FSB/BTG},{Alfa Inteligência},{Ideia},{Futura},{Gerp},{PoderData},{Veritá},{Ipespe},{AtlasIntel},{Quaest}}
\newcommand{\nPesos}{13}


\definecolor{azul}{HTML}{1F3A5F}
\definecolor{azulclaro}{HTML}{E8EEF5}
\definecolor{cinza}{HTML}{555555}
\definecolor{cinzaclaro}{HTML}{B5B5B5}
\definecolor{vermelho}{HTML}{B03A2E}
\definecolor{verde}{HTML}{2E7D4F}
\definecolor{codebg}{HTML}{F6F8FA}

\titleformat{\section}{\normalfont\Large\bfseries\color{azul}}{\thesection}{0.6em}{}
\titleformat{\subsection}{\normalfont\large\bfseries\color{azul!85!black}}{\thesubsection}{0.5em}{}
\titlespacing*{\section}{0pt}{14pt}{6pt}
\titlespacing*{\subsection}{0pt}{10pt}{4pt}
\setlength{\parindent}{0pt}
\setlength{\parskip}{5pt}
\setlist{itemsep=2pt,topsep=3pt,leftmargin=1.4em}

\pagestyle{fancy}
\fancyhf{}
\renewcommand{\headrulewidth}{0pt}
\fancyhead[R]{\footnotesize\color{cinza}Relatório técnico: previsão do 2º turno de 2026}
\fancyfoot[L]{\footnotesize\color{cinza}Valuche, Possamai e Lyra}
\fancyfoot[R]{\footnotesize\color{cinza}\thepage}

\lstdefinestyle{py}{
  language=Python, basicstyle=\ttfamily\footnotesize, backgroundcolor=\color{codebg},
  keywordstyle=\color{azul}\bfseries, commentstyle=\color{cinza}, stringstyle=\color{verde},
  showstringspaces=false, frame=single, rulecolor=\color{black!12}, framesep=4pt,
  xleftmargin=6pt, xrightmargin=6pt, breaklines=true, columns=fullflexible, keepspaces=true,
  literate={á}{{\'a}}1 {é}{{\'e}}1 {í}{{\'i}}1 {ó}{{\'o}}1 {ú}{{\'u}}1 {ã}{{\~a}}1 {õ}{{\~o}}1 {ç}{{\c{c}}}1 {Í}{{\'I}}1
}
\lstset{style=py}

\newtcolorbox{destaque}[1][]{colback=azulclaro,colframe=azul,boxrule=0.6pt,arc=2pt,left=6pt,right=6pt,top=4pt,bottom=4pt,#1}

\begin{document}

%================= CAPA =================
\begin{tcolorbox}[colback=azul,colframe=azul,arc=0pt,left=14pt,right=14pt,top=14pt,bottom=14pt]
\color{white}
{\small DESAFIO DE ESTATÍSTICA E ECONOMETRIA \textbullet\ FGV EPGE}\\[6pt]
{\huge\bfseries Previsão do 2º turno presidencial de 2026}\\[4pt]
{\Large Relatório técnico completo de modelagem econométrica}\\[10pt]
{\normalsize João Pedro Valuche de Andrade Pereira \textbullet\ Lethicia Manfioletti Possamai \textbullet\ Arthur Caron Lyra}
\end{tcolorbox}

\vspace{4pt}
\begin{destaque}
\textbf{Horário de corte dos dados: \horacorte{} (horário de Brasília).} Status da previsão: \textbf{\statusprevisao}.

Este relatório documenta integralmente a modelagem econométrica, a formulação analítica, o motor computacional, o backtest prospectivo e os testes de sensibilidade para a previsão do 2º turno da eleição presidencial brasileira de 2026 (Flávio Bolsonaro vs. Luiz Inácio Lula da Silva). Todos os números, gráficos e tabelas são reproduzidos de forma determinística pelo pipeline automatizado do repositório.
\end{destaque}

\vspace{6pt}
{\setlength{\tabcolsep}{0pt}
\begin{tabularx}{\linewidth}{@{}*{5}{>{\centering\arraybackslash}X}@{}}
\textcolor{azul}{\huge\bfseries \pFlavio\%} & \textcolor{azul}{\huge\bfseries \pLula\%} & \textcolor{cinza}{\huge\bfseries \pAbst\%} & \textcolor{cinza}{\huge\bfseries \pBrancos\%} & \textcolor{cinza}{\huge\bfseries \pNulos\%}\\[2pt]
\small Flávio Bolsonaro (PL) & \small Lula (PT) & \small Abstenção & \small Brancos & \small Nulos\\[-1pt]
\multicolumn{2}{c}{\footnotesize\color{cinza}\% dos votos válidos (Aba 1)} & \multicolumn{3}{c}{\footnotesize\color{cinza}\% do eleitorado apto (abstenção) e do comparecimento (brancos e nulos) (Aba 2)}
\end{tabularx}}

\vspace{8pt}
\setcounter{tocdepth}{2}
{\small\tableofcontents}

\newpage
%================= SEÇÃO 1 =================
\section{O problema do 2º turno e propriedades analíticas da métrica}

O Desafio de Estatística e Econometria da FGV EPGE exige duas previsões pontuais para o 2º turno presidencial de 25 de outubro de 2026:
\begin{enumerate}
    \item \textbf{Aba 1 (\texttt{Candidatos}):} Votos válidos (\%) para os dois candidatos finalistas, somando exatamente $100{,}0\%$ com uma casa decimal: Flávio Bolsonaro (PL) e Luiz Inácio Lula da Silva (PT).
    \item \textbf{Aba 2 (\texttt{Adicionais}):} Três agregados com uma casa decimal: a abstenção (como percentual do eleitorado apto total) e os votos brancos e nulos (como percentual do comparecimento às urnas).
\end{enumerate}

\subsection{Identidades da função objetivo no 2º turno \texorpdfstring{($K = 2$)}{(K = 2)}}

No 1º turno, com $K=12$ candidatos, a função perda dependia criticamente do tratamento de candidatos pequenos e da dispersão entre vários nomes. No 2º turno, como $K=2$ e os votos válidos somam $100\%$, ocorrem identidades matemáticas fundamentais que simplificam a estrutura da perda:

\begin{equation}
\hat{V}_{\text{Flávio}} + \hat{V}_{\text{Lula}} = 100\%, \qquad V_{\text{Flávio}} + V_{\text{Lula}} = 100\%.
\end{equation}

Disso decorre que o erro cometido em relação a um candidato é identicamente o erro do adversário com sinal oposto:
\begin{equation}
\hat{V}_{\text{Flávio}} - V_{\text{Flávio}} = - (\hat{V}_{\text{Lula}} - V_{\text{Lula}}).
\end{equation}

Assim, o Erro Absoluto Médio (MAE) da previsão de votos válidos reduz-se exatamente ao erro absoluto de um único candidato e a exatamente metade do erro na margem bipartidária:
\begin{equation}
\text{MAE}_{2T} = \frac{1}{2} \sum_{k=1}^{2} |\hat{V}_k - V_k| = |\hat{V}_{\text{Lula}} - V_{\text{Lula}}| = |\hat{V}_{\text{Flávio}} - V_{\text{Flávio}}| = \frac{1}{2} \left| (\hat{V}_{\text{Flávio}} - \hat{V}_{\text{Lula}}) - (V_{\text{Flávio}} - V_{\text{Lula}}) \right|.
\label{eq:mae_identidade}
\end{equation}

Essa propriedade implica que:
\begin{itemize}
    \item Errar a margem entre os finalistas em 2,0 pontos percentuais custa exatamente 1,0 ponto de MAE.
    \item Não há distorções provocadas por candidatos nanicos ou transferências parciais de terceiros na métrica de avaliação final: toda a estimação econométrica se resume a estimar com a maior acurácia possível uma única proporção escalar, $p \in [0, 1]$.
\end{itemize}

%================= SEÇÃO 2 =================
\section{Base de dados e auditoria}

\subsection{Resultados oficiais do 1º turno de 2026}
Diferentemente do 1º turno, no 2º turno dispõe-se de uma âncora empírica extraordinária: a apuração oficial do TSE do 1º turno, realizado 21 dias antes, em 04/10/2026. A base municipal e de zonas eleitorais do TSE foi processada com validação criptográfica (hash SHA-256):
\begin{itemize}
    \item \textbf{Votação válida no 1º turno:} Flávio Bolsonaro obteve \flavioUmT\% dos votos válidos, Lula obteve \lulaUmT\%, e os 10 candidatos eliminados somaram conjuntamente \tercUmT\% dos votos válidos. Na soma restrita aos dois primeiros, Flávio teve \flavioDoisCand\% e Lula teve \lulaDoisCand\%.
    \item \textbf{Comparecimento e agregados no 1º turno:} Abstenção de \abstUmT\% (sobre 158.745.079 eleitores aptos), votos brancos de \brancosUmT\% e votos nulos de \nulosUmT\% (sobre 125.275.835 votos totais apurados na urna).
\end{itemize}

\subsection{Séries históricas de 2º turno (2006 a 2022)}
Compilamos os microdados oficiais do TSE para todas as cinco eleições presidenciais com 2º turno na história recente (2006, 2010, 2014, 2018 e 2022). A base histórica contém:
\begin{itemize}
    \item Os percentuais oficiais de votos válidos dos finalistas, abstenção, brancos e nulos no 1º e no 2º turno, em nível nacional e por Unidade da Federação.
    \item \nHist{} pesquisas nacionais de intenção de voto estimuladas de 2º turno divulgadas entre o 1º e o 2º turno daquelas eleições pelos principais institutos do país.
    \item A relação das eleições estaduais para Governador com disputa de 2º turno ativa em cada ano.
\end{itemize}

\subsection{Pesquisas eleitorais do 2º turno de 2026}
Foram catalogadas \nPesq{} pesquisas estimuladas de 2º turno divulgadas por \nInst{} institutos de opinião pública entre 05/10/2026 e o horário de corte de 24/10/2026 às 20h00 BRT. Todas as pesquisas possuem registro auditável no sistema PesqEle do Tribunal Superior Eleitoral (TSE) e URLs canônicas verificadas.

Cada pesquisa é convertida em votos válidos descartando os votos brancos, nulos e indecisos:
\begin{equation}
\hat{v}_{j, \text{Lula}} = \frac{P_{j, \text{Lula}}}{P_{j, \text{Lula}} + P_{j, \text{Flávio}}}, \qquad \hat{v}_{j, \text{Flávio}} = 100\% - \hat{v}_{j, \text{Lula}}.
\end{equation}

Essa conversão distribui os eleitores indecisos proporcionalmente entre os decididos, hipótese operacional padrão equivalente à apuração de votos válidos pelo TSE.

%================= SEÇÃO 3 =================
\section{Pipeline e arquitetura computacional}

A Figura~\ref{fig:pipe2t} apresenta o fluxo completo de processamento de dados e geração dos entregáveis.

\begin{figure}[h]
\centering
\begin{tikzpicture}[node distance=0.45cm,
  caixa/.style={rectangle,rounded corners=3pt,draw=azul,fill=azul!8,text=azul!90!black,align=center,font=\scriptsize\bfseries,inner sep=4pt,minimum height=1.1cm},
  seta/.style={-{Stealth[length=5pt]},thick,azul},
  faixa/.style={rectangle,rounded corners=2pt,draw=vermelho!80,fill=vermelho!6,text=vermelho!80!black,align=center,font=\scriptsize\bfseries,inner sep=4pt}]
\node[caixa,text width=2.4cm] (a) {TSE 2026 (1º Turno)\\+ PesqEle 2º Turno};
\node[caixa,text width=2.1cm,right=of a] (b) {Conversão em\\votos válidos};
\node[caixa,text width=2.3cm,right=of b] (c) {Última pesquisa\\por instituto};
\node[caixa,text width=2.4cm,right=of c] (d) {Ponderação M2\\$w_j = \frac{1}{e_{j,1T}^2 + \sigma_0^2}$};
\node[caixa,text width=2.0cm,right=of d] (e) {Aba 2 Ancorada\\Zero\_Delta};
\node[caixa,text width=1.8cm,right=of e] (f) {Maiores\\restos};
\foreach \x/\y in {a/b,b/c,c/d,d/e,e/f} \draw[seta] (\x)--(\y);
\node[faixa,below=0.45cm of c,text width=8.5cm] (v) {Protocolo pré-registrado de validação prospectiva em Expanding Window\\(2014 treinado em 2006--2010; 2018 em 2006--2014; 2022 em 2006--2018)};
\draw[seta,vermelho] (v.north -| d) -- (d);
\end{tikzpicture}
\caption{Diagrama do pipeline do 2º turno. Em vermelho, a etapa de validação temporal.}
\label{fig:pipe2t}
\end{figure}

O código a seguir exemplifica o cálculo dos pesos $w_j$ no modelo M2 com base no erro observado no 1º turno:

\begin{lstlisting}
def estimar_m2_2t(df_pesquisas_2t, erros_1t_map, sigma0=1.5):
    """
    Agregador M2: pondera a pesquisa de 2o turno de cada instituto pelo inverso
    do erro quadratico cometido no 1o turno da mesma eleicao, regularizado por sigma0^2.
    """
    ultimas = (df_pesquisas_2t.sort_values("data_divulgacao", ascending=False)
                              .drop_duplicates(subset=["instituto_clean"]))
    mediana_erro_1t = np.median([abs(e) for e in erros_1t_map.values()])
    
    pesos, valores_pt = [], []
    for _, linha in ultimas.iterrows():
        inst = linha["instituto_clean"]
        erro_1t = erros_1t_map.get(inst, mediana_erro_1t)
        peso = 1.0 / (erro_1t**2 + sigma0**2)
        v_pt = linha["pct_pt_validos"]
        pesos.append(peso)
        valores_pt.append(v_pt)
        
    p_pt_estimado = np.average(valores_pt, weights=pesos)
    return p_pt_estimado, 100.0 - p_pt_estimado
\end{lstlisting}

%================= SEÇÃO 4 =================
\section{Modelos candidatos e formulação econométrica}

\subsection{Etapa 1: Agregadores base}

\paragraph{M0\_2T (Média Simples):}
Média aritmética simples da última pesquisa de 2º turno de cada instituto:
\begin{equation}
\hat{p}_{\text{PT}}^{\text{M0}} = \frac{1}{J} \sum_{j=1}^{J} \hat{v}_{j, \text{PT}}.
\end{equation}
Não possui hiperparâmetros livres e serve como referência estrita de parcimônia.

\paragraph{M1\_2T (Recência e Amostra):}
Pondera todas as pesquisas divulgadas na janela do 2º turno pelo tamanho da amostra e por decaimento exponencial temporal com meia-vida $h \in \{3, 5, 7\}$ dias:
\begin{equation}
w_i = \sqrt{N_i} \cdot 2^{-\Delta t_i / h}, \qquad \hat{p}_{\text{PT}}^{\text{M1}} = \frac{\sum_i w_i \hat{v}_{i, \text{PT}}}{\sum_i w_i}.
\end{equation}

\paragraph{M2\_2T (Ponderação pela Precisão Empírica no 1º Turno):}
Aproveita a observabilidade do erro amostral intra-ciclo. Definimos o semi-erro da margem entre os dois primeiros colocados no 1º turno:
\begin{equation}
e_{j, 1T} = \frac{1}{2}\left[ (\hat{v}_{j, \text{PT}}^{\,1T} - \hat{v}_{j, \text{Adv}}^{\,1T}) - (v_{\text{PT}}^{\,1T} - v_{\text{Adv}}^{\,1T}) \right].
\end{equation}
O peso de cada instituto no 2º turno é dado por:
\begin{equation}
w_j = \frac{1}{e_{j, 1T}^2 + \sigma_0^2}, \qquad \hat{p}_{\text{PT}}^{\text{M2}} = \frac{\sum_j w_j \hat{v}_{j, \text{PT}}}{\sum_j w_j}, \qquad \sigma_0 \in \{1{,}5; \, 2{,}5\}\text{ p.p.}
\end{equation}
O parâmetro de regularização $\sigma_0$ impede que institutos que acertaram a urna por pura variância amostral recebam peso infinito. Institutos sem pesquisa de véspera no 1º turno recebem a mediana do erro absoluto dos demais institutos daquela eleição.

\subsection{Etapa 2: Ajustes avaliados sobre o agregador base}

\paragraph{Ajuste 1 --- Persistência parcial intra-ciclo ($\rho$):}
Subtrai de cada pesquisa de 2º turno uma fração $\rho$ do erro cometido pelo mesmo instituto no 1º turno:
\begin{equation}
\hat{v}_{j, \text{PT}}^{\text{corr}} = \hat{v}_{j, \text{PT}} - \rho \cdot e_{j, 1T}, \qquad \rho \in \{0{,}20; \, 0{,}35; \, \hat{\rho}_{\text{OLS}}\}.
\end{equation}

\paragraph{Ajuste 2 --- Viés comum histórico com shrinkage ($k_\mu$):}
Desconta o erro médio histórico cometido pelos institutos nos 2ºs turnos passados com amortecimento empírico:
\begin{equation}
\hat{\mu}_{\text{PT}, 2T} = \frac{N_{\text{treino}}}{N_{\text{treino}} + k_\mu} \cdot \bar{e}_{\text{PT}, 2T}^{\text{treino}}, \qquad k_\mu \in \{1, 3\}.
\end{equation}

\paragraph{Ajuste 3 --- Shrinkage com prior estrutural de transferência do 1º turno ($\lambda$):}
Combina a estimativa das pesquisas com a projeção estrutural decorrente da matriz de transferência dos votos do 1º turno:
\begin{equation}
\hat{p}_{\text{PT}} = (1 - \lambda) \hat{p}_{\text{PT}}^{\text{pesq}} + \lambda \, p_{\text{PT}}^{\text{prior}}, \qquad \lambda \in \{0{,}10; \, 0{,}25\}.
\end{equation}

%================= SEÇÃO 5 =================
\section{Protocolo de validação prospectiva e resultados}

A regra formal pré-registrada exige que qualquer modelo ou ajuste candidato reduza o MAE médio da janela crescente (\emph{Expanding Window}) em mais de um erro-padrão da diferença em relação ao modelo base:
\begin{equation}
\bar{\Delta} = \text{MAE}_{\text{cand}} - \text{MAE}_{\text{base}} < -\text{SE}(\Delta).
\end{equation}
Caso contrário, o modelo mais simples é retido por parcimônia.

\subsection{Resultados da Etapa 1 e Etapa 2}

A Tabela~\ref{tab:modelos} reporta o desempenho prospectivo de todos os modelos e ajustes testados.

\begin{table}[h]
\centering\small
\caption{Validação temporal do 2º turno em Expanding Window (MAE em p.p.).}
\label{tab:modelos}
\begin{tabular}{lcccccl}
\toprule
\textbf{Etapa 1: agregador} & \textbf{2014} & \textbf{2018} & \textbf{2022} & \textbf{Média} & \textbf{$\bar\Delta$ / EP} & \textbf{Decisão} \\
\midrule
M0 média simples & 0,63 & 1,82 & 0,29 & 0,92 &  & referência \\
M1 recência e amostra ($h$=3) & 0,77 & 0,96 & 0,45 & 0,72 & +0,19 / 0,34 & não supera M0 \\
M2 peso pelo erro no 1º turno ($\sigma_0$=1,5) & 0,56 & 1,29 & 0,10 & 0,65 & +0,27 / 0,14 & \textbf{escolhido} \\
M2 peso pelo erro no 1º turno ($\sigma_0$=2,5) & 0,56 & 1,27 & 0,11 & 0,65 & +0,27 / 0,15 & supera, MAE maior \\
\midrule
\textbf{Etapa 2: ajustes sobre o M2} & & & & & & \\
\midrule
+ correção pelo erro do 1º turno ($\hat\rho$ MQO) & 0,29 & 1,40 & 0,01 & 0,57 & +0,08 / 0,11 & empate: fica a base \\
+ correção pelo erro do 1º turno ($\rho=0{,}2$) & 0,51 & 1,44 & 0,34 & 0,77 & $-$0,12 / 0,09 & piora \\
+ viés comum histórico do 2º turno ($k_\mu=3$) & 0,48 & 1,45 & 0,05 & 0,66 & $-$0,01 / 0,08 & empate: fica a base \\
+ 10\% do prior de transferência ($\tilde\tau$) & 1,61 & 1,62 & 0,12 & 1,12 & $-$0,47 / 0,31 & piora \\
\bottomrule
\end{tabular}

\end{table}

\textbf{Conclusões centrais da validação:}
\begin{itemize}
    \item O modelo \textbf{M2\_2T ($\sigma_0=1{,}5$)} foi o vencedor inequívoco da Etapa 1. Ele reduziu o erro médio de 0,92 p.p. (M0) para 0,65 p.p., superando a média simples por expressivos \textbf{1,95 erro-padrão} ($\bar{\Delta} = -0{,}27$, $\text{EP} = 0{,}14$).
    \item Na validação complementar LOEO (Tabela~\ref{tab:loeo}), o M2 superou a média simples em \emph{todas} as cinco eleições avaliadas (2006, 2010, 2014, 2018 e 2022), obtendo MAE médio global de 0,48 p.p. contra 0,77 p.p. do M0.
    \item Nenhum ajuste da Etapa 2 foi aprovado. A persistência intra-ciclo ($\rho$) piorou ou empatou tecnicamente com o modelo puro, demonstrando que os institutos tendem a recalibrar suas ponderações e cotas amostrais entre os turnos, evitando a repetição estrita do erro pontual do 1º turno.
    \item O modelo final oficial selecionado para a Aba 1 é, portanto, o \textbf{M2\_2T ($\sigma_0=1{,}5$) Puro}.
\end{itemize}

\begin{table}[h]
\centering\small
\caption{Validação complementar Leave-One-Election-Out (LOEO, 2006 a 2022).}
\label{tab:loeo}
\begin{tabular}{lcccccc}
\toprule
\textbf{Modelo} & \textbf{2006} & \textbf{2010} & \textbf{2014} & \textbf{2018} & \textbf{2022} & \textbf{Média LOEO (EP)} \\
\midrule
M0\_2T & 0,88 & 0,25 & 0,63 & 1,82 & 0,29 & 0,77 (0,29) \\
M1\_2T (h=3) & 0,52 & 0,17 & 0,77 & 0,96 & 0,45 & 0,57 (0,14) \\
M1\_2T (h=5) & 0,47 & 0,22 & 0,49 & 1,42 & 0,42 & 0,60 (0,21) \\
M1\_2T (h=7) & 0,44 & 0,24 & 0,31 & 1,67 & 0,41 & 0,61 (0,27) \\
\textbf{M2\_2T (sigma0=1.5)} & 0,44 & 0,03 & 0,56 & 1,29 & 0,10 & \textbf{0,48} (0,23) \\
M2\_2T (sigma0=2.5) & 0,46 & 0,02 & 0,56 & 1,27 & 0,11 & 0,49 (0,22) \\
M0\_2T + Persist 1T (rho=0.20) & 0,15 & 0,75 & 0,46 & 2,18 & 0,50 & 0,81 (0,36) \\
M0\_2T + Persist 1T (rho=0.35) & 0,93 & 1,49 & 1,27 & 2,45 & 0,65 & 1,36 (0,31) \\
M2\_2T (sigma0=1.5) + Persist 1T (rho=0.20) & 0,29 & 0,96 & 0,51 & 1,44 & 0,34 & 0,71 (0,22) \\
M2\_2T (sigma0=1.5) + Persist 1T (rho=0.35) & 0,84 & 1,65 & 1,32 & 1,56 & 0,52 & 1,18 (0,22) \\
\bottomrule
\end{tabular}

\end{table}

\subsection{Pesos atribuídos a cada instituto em 2026}

A Tabela~\ref{tab:pesos} detalha o semi-erro na margem do 1º turno e o peso resultante de cada instituto na agregação final do 2º turno de 2026.

\begin{table}[h]
\centering\small
\caption{Calibração dos pesos no modelo M2\_2T para as pesquisas de 2026.}
\label{tab:pesos}
\begin{tabular}{lcccc}
\toprule
\textbf{Instituto} & \textbf{Semi-Erro 1T ($e_{j,1T}$)} & \textbf{Origem do Erro} & \textbf{Peso Absoluto ($w_j$)} & \textbf{Participação Relativa (\%)} \\
\midrule
Quaest & 1,43 p.p. & Pesquisa de véspera 1T & 0,232 & 13,4\% \\
AtlasIntel & 2,38 p.p. & Pesquisa de véspera 1T & 0,126 & 7,3\% \\
Alfa Inteligência & 2,42 p.p. & Mediana histórica & 0,123 & 7,1\% \\
Ipespe & 2,42 p.p. & Mediana histórica & 0,123 & 7,1\% \\
Veritá & 2,42 p.p. & Mediana histórica & 0,123 & 7,1\% \\
FSB/BTG & 2,42 p.p. & Mediana histórica & 0,123 & 7,1\% \\
Palver & 2,42 p.p. & Mediana histórica & 0,123 & 7,1\% \\
Gerp & 2,42 p.p. & Mediana histórica & 0,123 & 7,1\% \\
Ideia & 2,42 p.p. & Mediana histórica & 0,123 & 7,1\% \\
PoderData & 2,42 p.p. & Mediana histórica & 0,123 & 7,1\% \\
Futura & 2,42 p.p. & Mediana histórica & 0,123 & 7,1\% \\
Datafolha & 2,46 p.p. & Pesquisa de véspera 1T & 0,120 & 6,9\% \\
Real Time Big Data & 3,08 p.p. & Pesquisa de véspera 1T & 0,085 & 4,9\% \\
CNT/MDA & 3,78 p.p. & Pesquisa de véspera 1T & 0,060 & 3,5\% \\
\bottomrule
\end{tabular}

\end{table}

%================= SEÇÃO 6 =================
\section{Transferência de votos e checagem independente de plausibilidade}

O resultado oficial do 1º turno permite realizar uma verificação estrutural independente da previsão das pesquisas. Definimos $\tau_t$ como a fração dos votos dos candidatos eliminados no 1º turno que migrou para o candidato do PT no 2º turno:
\begin{equation}
\tau_t = \frac{V_{\text{PT}, 2T, t} - V_{\text{PT}, 1T, t}}{100 - V_{\text{PT}, 1T, t} - V_{\text{Adv}, 1T, t}}.
\end{equation}

A Tabela~\ref{tab:transf} apresenta a série histórica de transferência em todas as eleições com 2º turno.

\begin{table}[h]
\centering\small
\caption{Transferência aparente de votos do 1º para o 2º turno (\% dos válidos).}
\label{tab:transf}
\begin{tabular}{lcccccc}
\toprule
 & 2006 & 2010 & 2014 & 2018 & 2022 & \textbf{2026} \\
\midrule
PT no 1º turno & 48,6 & 46,9 & 41,6 & 29,3 & 48,4 & 45,2 \\
Demais candidatos no 1º turno & 9,8 & 20,5 & 24,9 & 24,7 & 8,4 & 7,8 \\
PT no 2º turno & 60,8 & 56,1 & 51,6 & 44,9 & 50,9 & ? \\
$\tau$ (fração capturada pelo PT) & 1,25 & 0,45 & 0,40 & 0,63 & 0,30 & $\tilde\tau=0,45$ \\
\bottomrule
\end{tabular}

\end{table}

\subsection{O outlier de 2006 e a superioridade da mediana}
A eleição de 2006 apresenta $\tau_{2006} = 1{,}2530 > 1$. Isso ocorreu porque Geraldo Alckmin (PSDB) teve uma queda anômala de votação entre os turnos: obteve 41,64\% no 1º turno e caiu para 39,17\% no 2º turno (perda nominal de mais de 2 milhões de votos). Nesse caso, Lula não capturou apenas votos de terceiros candidatos, mas absorveu eleitores que haviam votado no próprio adversário no 1º turno.

Por essa razão, a média histórica ($\bar{\tau} = 0{,}61$) é severamente inflada. O estimador estatisticamente robusto para a transferência típica é a \textbf{mediana histórica}, $\tilde{\tau} = \mathbf{0{,}45}$ (44,6\%).

\subsection{Convergência estrutural com a previsão das pesquisas}
Aplicando a taxa mediana histórica de transferência ($\tilde{\tau} = 0{,}4464$) sobre os votos dos candidatos eliminados no 1º turno de 2026 (7,81\% dos votos válidos):
\begin{equation}
\hat{V}_{\text{Lula}}^{\text{prior}} = 45{,}16\% + 0{,}4464 \times 7{,}81\% = 45{,}16\% + 3{,}49\% = \mathbf{48{,}65\%} \implies \mathbf{48{,}6\%}.
\end{equation}
\begin{equation}
\hat{V}_{\text{Flávio}}^{\text{prior}} = 100\% - 48{,}65\% = \mathbf{51{,}35\%} \implies \mathbf{51{,}4\%}.
\end{equation}

Essa verificação independente corrobora com impressionante precisão a estimativa agregada das pesquisas (\pFlavio\% para Flávio Bolsonaro vs. \pLula\% para Lula), confirmando a robustez fundamental do modelo.

%================= SEÇÃO 7 =================
\section{Modelagem da Aba 2 (Abstenção, brancos e nulos)}

No 1º turno, a projeção da abstenção por tendência linear histórica superestimou a abstenção real em 1,22 p.p. No 2º turno, a presença do resultado oficial do 1º turno ancorou cinco famílias de modelos prospectivos no Expanding Window (Tabela~\ref{tab:aba2}):
\begin{itemize}
    \item \texttt{Zero\_Delta}: $\hat{y}_{2T} = y_{1T}$ (parcimônia estrita: repetir o valor observado 21 dias antes).
    \item \texttt{Persistencia\_Delta}: $\hat{y}_{2T} = y_{1T} + \Delta y_{t-1}$.
    \item \texttt{Media\_Delta}: $\hat{y}_{2T} = y_{1T} + \bar{\Delta}_{\text{treino}}$.
    \item \texttt{Ratio\_Media}: $\hat{y}_{2T} = y_{1T} \times \bar{r}_{\text{treino}}$.
    \item \texttt{BottomUp\_UF\_Gov2T}: Projeção estado por estado discriminando a presença de 2º turno estadual para Governador.
\end{itemize}

\begin{table}[h]
\centering\small
\caption{Backtest da Aba 2 no Expanding Window (2014, 2018, 2022) e previsão final.}
\label{tab:aba2}
\begin{tabular}{lccc}
\toprule
\textbf{Método} & \textbf{Abstenção} & \textbf{Brancos} & \textbf{Nulos} \\
\midrule
Repetir o 1º turno & \textbf{1,01} & \textbf{0,93} & \textbf{0,93} \\
1º turno + variação da eleição anterior & 1,25 & 1,09 & 1,15 \\
1º turno + variação média & 1,67 & 1,01 & 1,11 \\
1º turno $\times$ razão média 2T/1T & 1,98 & 0,58 & 1,05 \\
Por UF, com e sem 2º turno estadual & 1,38 & n.a. & n.a. \\
\midrule
1º turno de 2026 (TSE) & 21,08 & 1,84 & 2,93 \\
\textbf{Previsão para o 2º turno} & \textbf{21,1} & \textbf{1,8} & \textbf{2,9} \\
\bottomrule
\end{tabular}

\end{table}

Pela regra pré-registrada de parcimônia, nenhum modelo concorrente superou \texttt{Zero\_Delta} por mais de 1 erro-padrão. Assim, a previsão oficial para a Aba 2 é a repetição dos percentuais apurados no 1º turno de 2026:
\textbf{Abstenção: \pAbst\%}, \textbf{Votos Brancos: \pBrancos\%} e \textbf{Votos Nulos: \pNulos\%}.

\subsection{Limitação estrutural documentada: 2º turno para Governador}
Em 2026, apenas 6 Unidades da Federação (AC, AM, DF, ES, RN e TO) realizam 2º turno para Governador. Esses estados representam somente \textbf{\govShareVinteSeis\%} do eleitorado nacional apto. No período 2006--2022, a média histórica de eleitores com disputa estadual ativa de 2º turno foi de \textbf{\govShareHist\%}. 

Historicamente, nos estados sem disputa estadual, a abstenção subiu em média $+0{,}12$ p.p. a mais do que nos estados com disputa. Embora o modelo de decomposição não tenha superado o baseline por mais de um erro-padrão no backtest ($n=3$), a quase ausência de campanhas estaduais ativas no 2º turno de 2026 constitui uma limitação estrutural documentada.

%================= SEÇÃO 8 =================
\section{Análise de sensibilidade}

A Tabela~\ref{tab:sens} compara a previsão oficial com todas as especificações alternativas avaliadas.

\begin{table}[h]
\centering\small
\caption{Grade de sensibilidade para a eleição presidencial de 2026 (2º turno).}
\label{tab:sens}
\begin{tabular}{lccc}
\toprule
\textbf{Especificação / Modelo Alternativo} & \textbf{Flávio Cont.} & \textbf{Flávio (1 casa)} & \textbf{Lula (1 casa)} \\
\midrule
M0\_2T (Média Simples Última de cada Instituto) & 51,32\% & 51,3\% & 48,7\% \\
M1\_2T (h=3 dias, Recência + Raiz Amostra) & 51,09\% & 51,1\% & 48,9\% \\
M1\_2T (h=5 dias, Recência + Raiz Amostra) & 51,16\% & 51,2\% & 48,8\% \\
M1\_2T (h=7 dias, Recência + Raiz Amostra) & 51,20\% & 51,2\% & 48,8\% \\
M2\_2T (sigma0=1.5, Peso Inverso Erro 1T) & 51,29\% & 51,3\% & 48,7\% \\
M2\_2T (sigma0=2.5, Peso Inverso Erro 1T) & 51,31\% & 51,3\% & 48,7\% \\
M0\_2T + Persistência 1T (rho=0.20) & 51,83\% & 51,8\% & 48,2\% \\
M2\_2T (sigma0=1.5) + Persistência 1T (rho=0.20) & 51,78\% & 51,8\% & 48,2\% \\
M2\_2T (sigma0=1.5) + Persistência 1T (rho\_OLS=0.15) & 51,67\% & 51,7\% & 48,3\% \\
Prior transferência com tau médio (tau=0.6061) & 50,10\% & 50,1\% & 49,9\% \\
Prior transferência com tau mediano (tau=0.4464) & 51,35\% & 51,4\% & 48,6\% \\
\textbf{\textbf{Modelo Oficial: M2\_2T (sigma0=1.5) (Puro)}} & \textbf{51,29\%} & \textbf{51,3\%} & \textbf{48,7\%} \\
\bottomrule
\end{tabular}

\end{table}

Em todas as especificações plausíveis, a liderança numérica de Flávio Bolsonaro oscila em uma faixa muito estreita, entre 51,0\% e 51,8\%, com Lula entre 48,2\% e 49,0\%, atestando a estabilidade da previsão frente a escolhas metodológicas alternativas.

\newpage
%================= APÊNDICE =================
\appendix
\section{Relação completa de pesquisas de 2º turno catalogadas no TSE}

A Tabela~\ref{tab:todas_pesquisas} lista todas as pesquisas nacionais de 2º turno cadastradas em nossa base de dados até o corte de \horacorte.

{\scriptsize\setlength{\tabcolsep}{3.5pt}
\begin{longtable}{llccccccc}
\caption{Pesquisas nacionais de 2º turno da eleição presidencial de 2026 registradas no TSE.}\label{tab:todas_pesquisas}\\
\toprule
\textbf{Instituto} & \textbf{Registro TSE} & \textbf{Campo} & \textbf{$N$} & \textbf{Método} & \textbf{Flávio} & \textbf{Lula} & \textbf{Vál. Flávio} & \textbf{Vál. Lula} \\
\midrule
\endfirsthead
\multicolumn{9}{c}{\footnotesize\color{cinza}\tablename\ \thetable\ -- continuação da página anterior}\\
\toprule
\textbf{Instituto} & \textbf{Registro TSE} & \textbf{Campo} & \textbf{$N$} & \textbf{Método} & \textbf{Flávio} & \textbf{Lula} & \textbf{Vál. Flávio} & \textbf{Vál. Lula} \\
\midrule
\endhead
\midrule
\multicolumn{9}{r}{\footnotesize\color{cinza}continua na próxima página\dots}\\
\endfoot
\bottomrule
\endlastfoot
FSB/BTG & BR-08234/2026 & 23/10--24/10 & 2.000 & Telefonica & 47,0 & 46,0 & 50,5\% & 49,5\% \\
Quaest & BR-09158/2026 & 23/10--24/10 & 3.702 & Presencial & 51,0 & 49,0 & 51,0\% & 49,0\% \\
Datafolha & BR-09540/2026 & 23/10--24/10 & 4.050 & Presencial & 48,0 & 47,0 & 51,0\% & 49,0\% \\
CNT/MDA & BR-02136/2026 & 21/10--23/10 & 2.032 & Presencial & 46,3 & 45,0 & 50,7\% & 49,3\% \\
Real Time Big Data & BR-09670/2026 & 19/10--22/10 & 2.000 & Telefonica & 49,0 & 45,0 & 52,0\% & 48,0\% \\
Veritá & BR-06070/2026 & 18/10--22/10 & 15.554 & Telefonica & 50,2 & 47,5 & 51,3\% & 48,7\% \\
Alfa Inteligência & BR-08496/2026 & 18/10--22/10 & 2.700 & Telefonica & 49,9 & 46,0 & 52,0\% & 48,0\% \\
Ipespe & BR-04350/2026 & 20/10--22/10 & 1.500 & Telefonica & 50,5 & 49,5 & 50,5\% & 49,5\% \\
Datafolha & BR-09709/2026 & 20/10--22/10 & 2.520 & Presencial & 47,0 & 46,0 & 50,0\% & 50,0\% \\
PoderData & BR-07926/2026 & 19/10--21/10 & 4.000 & Telefonica & 49,0 & 48,0 & 50,5\% & 49,5\% \\
Quaest & BR-08535/2026 & 17/10--20/10 & 2.154 & Presencial & 45,0 & 44,0 & 50,6\% & 49,4\% \\
Ideia & BR-08251/2026 & 16/10--20/10 & 1.500 & Telefonica & 49,3 & 47,2 & 51,1\% & 48,9\% \\
Futura & BR-06614/2026 & 19/10--20/10 & 2.000 & Telefonica & 48,9 & 46,8 & 51,1\% & 48,9\% \\
Palver & BR-04612/2026 & 18/10--20/10 & 5.000 & Online & 50,0 & 48,0 & 51,0\% & 49,0\% \\
Real Time Big Data & BR-02952/2026 & 17/10--19/10 & 2.000 & Telefonica & 48,0 & 45,0 & 52,0\% & 48,0\% \\
Alfa Inteligência & BR-00820/2026 & 15/10--19/10 & 2.700 & Telefonica & 48,8 & 45,9 & 51,5\% & 48,5\% \\
AtlasIntel & BR-02418/2026 & 13/10--18/10 & 5.000 & Online & 49,8 & 47,6 & 51,1\% & 48,9\% \\
PoderData & BR-02586/2026 & 16/10--18/10 & 3.000 & Telefonica & 49,0 & 46,0 & 52,0\% & 48,0\% \\
FSB/BTG & BR-04499/2026 & 16/10--18/10 & 2.000 & Telefonica & 46,5 & 45,9 & 50,3\% & 49,7\% \\
Futura & BR-06715/2026 & 16/10--17/10 & 2.000 & Telefonica & 48,3 & 47,5 & 50,4\% & 49,6\% \\
Veritá & BR-01107/2026 & 12/10--16/10 & 15.554 & Telefonica & 48,4 & 46,3 & 51,1\% & 48,9\% \\
Ipespe & BR-00223/2026 & 13/10--15/10 & 1.500 & Telefonica & 51,2 & 48,8 & 51,2\% & 48,8\% \\
CNT/MDA & BR-00652/2026 & 13/10--15/10 & 2.002 & Presencial & 46,8 & 45,9 & 50,5\% & 49,5\% \\
Futura & BR-03764/2026 & 14/10--15/10 & 2.000 & Telefonica & 49,0 & 47,5 & 50,8\% & 49,2\% \\
PoderData & BR-04258/2026 & 13/10--15/10 & 3.000 & Telefonica & 47,0 & 47,0 & 50,0\% & 50,0\% \\
Datafolha & BR-00466/2026 & 13/10--15/10 & 2.520 & Presencial & 47,0 & 45,0 & 51,0\% & 49,0\% \\
Palver & BR-08028/2026 & 10/10--13/10 & 5.000 & Online & 50,0 & 49,0 & 50,5\% & 49,5\% \\
Quaest & BR-08729/2026 & 10/10--13/10 & 2.154 & Presencial & 46,0 & 44,0 & 51,1\% & 48,9\% \\
PoderData & BR-01209/2026 & 10/10--12/10 & 3.000 & Telefonica & 48,0 & 46,0 & 51,0\% & 49,0\% \\
FSB/BTG & BR-07997/2026 & 10/10--12/10 & 2.000 & Telefonica & 48,0 & 45,0 & 52,0\% & 48,0\% \\
Gerp & BR-07877/2026 & 08/10--12/10 & 2.400 & Presencial & 53,0 & 44,0 & 55,0\% & 45,0\% \\
Futura & BR-00650/2026 & 07/10--10/10 & 2.000 & Telefonica & 50,0 & 46,6 & 51,8\% & 48,2\% \\
Ipespe & BR-00933/2026 & 06/10--08/10 & 1.500 & Telefonica & 49,0 & 44,0 & 52,7\% & 47,3\% \\
AtlasIntel & BR-03663/2026 & 03/10--08/10 & 5.026 & Online & 51,1 & 45,7 & 52,8\% & 47,2\% \\
Datafolha & BR-02949/2026 & 06/10--07/10 & 2.520 & Presencial & 49,0 & 45,0 & 52,0\% & 48,0\% \\
PoderData & BR-08134/2026 & 05/10--07/10 & 3.000 & Telefonica & 49,0 & 44,0 & 53,0\% & 47,0\% \\
\end{longtable}
}

\end{document}
